\chapter[1989 --- Bishop et al.]{1989: J. Michael Bishop, Harold E. Varmus}
\label{ch:1989_Bishop}

\section*{Main Contribution}

\textit{Discovered that cancer-causing genes (oncogenes) arise from normal cellular genes (proto-oncogenes) that have been mutated, explaining the genetic origin of cancer.}

\section{Official Citation}

for their discovery that oncogenes arise from proto-oncogenes

The 1989 Nobel Prize in Physiology or Medicine was awarded jointly to J. Michael Bishop and Harold E. Varmus for their discovery that oncogenes---genes that can cause cancer when they are activated or mutated---arise from normal cellular genes called proto-oncogenes. This discovery fundamentally changed our understanding of cancer, revealing that cancer is not simply the result of foreign agents (such as viruses) attacking the body from the outside, but rather the consequence of the corruption of the body's own genetic machinery. The identification of proto-oncogenes as the precursors of oncogenes opened up entirely new approaches to the diagnosis, classification, and treatment of cancer.

\section{Background and Historical Context}

The study of cancer has been one of the most challenging and frustrating areas of medical research. For much of the twentieth century, cancer was understood primarily as a disease of uncontrolled cell growth: cancer cells divide without the normal constraints of growth regulation, forming tumors that can invade surrounding tissues and spread (metastasize) to distant organs. However, the molecular mechanisms that cause cells to become cancerous were unknown.

The study of tumor viruses---viruses that can cause cancer in animals---provided an important avenue of investigation. In the 1960s and 1970s, researchers discovered that certain RNA tumor viruses (retroviruses) could cause cancers in chickens, mice, and other animals. These retroviruses were found to carry genes---called viral oncogenes (v-onc)---that were responsible for their cancer-causing ability. When the sequences of viral oncogenes were compared with the sequences of normal cellular DNA, a surprising discovery was made: the viral oncogenes had close counterparts in the normal cellular genome. These normal cellular counterparts were called proto-oncogenes (c-onc), and they appeared to be ordinary cellular genes that had been ``captured'' and modified by retroviruses during their evolution.

J. Michael Bishop was born in York, Pennsylvania, in 1936, and trained in medicine at Harvard Medical School. He joined the faculty of the University of California, San Francisco (UCSF), in 1968, where he began studying the molecular biology of retroviruses. Harold Eliot Varmus was born in Oelwein, Iowa, in 1939, and trained in medicine at Columbia University and in English literature at Amherst College. He joined the UCSF faculty in 1972, where he began collaborating with Bishop on the study of retroviruses and cancer.

Bishop and Varmus's collaboration began in the early 1970s, when they decided to investigate the origin of viral oncogenes. The prevailing hypothesis was that viral oncogenes had been derived from normal cellular genes, but this hypothesis had not been directly tested. Bishop and Varmus set out to determine whether the oncogene of Rous sarcoma virus (RSV), called \emph{src}, was related to any normal cellular gene.

\section{The Discovery/Contribution}

\subsection{The Search for Cellular Oncogenes}

Bishop and Varmus's approach was to use molecular hybridization---a technique in which radioactively labeled DNA probes are used to detect complementary sequences in other DNA molecules---to search for cellular genes that were similar in sequence to the viral \emph{src} oncogene. In a series of experiments conducted between 1975 and 1976, they showed that DNA from normal, uninfected chicken cells contained sequences that were highly similar (but not identical) to the viral \emph{src} gene. This was a landmark finding: it demonstrated that the oncogene of Rous sarcoma virus was derived from a normal cellular gene---the proto-oncogene \emph{c-src}---that existed in the genome of every cell in the chicken's body.

Bishop and Varmus published their seminal paper in \emph{Cell} in 1976, and the discovery was immediately recognized as a major advance. It demonstrated that the distinction between oncogenes and proto-oncogenes was not a qualitative one (oncogenes were not fundamentally different genes from proto-oncogenes), but rather a quantitative and qualitative one: oncogenes were proto-oncogenes that had been activated, overexpressed, or mutated by retroviruses (or by other mechanisms) in ways that caused them to drive uncontrolled cell growth.

\subsection{The Diversity of Proto-oncogenes and Their Mechanisms of Activation}

Following the discovery that \emph{c-src} was the cellular counterpart of the viral \emph{src} oncogene, Bishop and Varmus and many other groups showed that virtually all known viral oncogenes had cellular counterparts. Each retroviral oncogene was found to be derived from a specific cellular proto-oncogene, and these proto-oncogenes encoded a diverse array of proteins involved in cell growth and proliferation, including growth factors, growth factor receptors, signal transduction molecules, transcription factors, and cell cycle regulators.

Bishop and Varmus also elucidated the mechanisms by which proto-oncogenes can be activated to become oncogenes. They showed that proto-oncogenes can be activated in several ways:

\begin{itemize}
 \item \textbf{Insertional mutagenesis:} When a retrovirus integrates its DNA near a proto-oncogene, the viral promoter or enhancer sequences can activate the proto-oncogene, causing it to be overexpressed.
 \item \textbf{Point mutation:} A single nucleotide change in a proto-oncogene can alter the structure or function of its protein product, converting it from a regulated growth signal to a constitutive (always-on) growth signal. For example, point mutations in the \emph{RAS} proto-oncogene are found in approximately 30\% of all human cancers.
 \item \textbf{Gene amplification:} Extra copies of a proto-oncogene can be produced by gene amplification, leading to overexpression of the encoded protein. Gene amplification of the \emph{MYC} proto-oncogene, for example, is a hallmark of many types of cancer.
 \item \textbf{Chromosomal translocation:} A chromosomal translocation can juxtapose a proto-oncogene with a highly expressed gene, causing the proto-oncogene to be expressed at abnormally high levels. The translocation t(9;22), which produces the BCR-ABL fusion protein in chronic myeloid leukemia, is a well-known example.
\end{itemize}

Bishop and Varmus's work thus provided a unifying framework for understanding the molecular basis of cancer: cancer is caused by the activation of proto-oncogenes (by retroviruses, mutations, amplifications, or translocations) and by the inactivation of tumor suppressor genes (genes that normally inhibit cell growth). This framework---the ``two-hit'' model of cancer, in which both activating mutations in oncogenes and inactivating mutations in tumor suppressor genes are required for cancer development---has become the foundation of modern cancer biology.

\subsection{The Discovery of Oncogenes in Human Cancer}

One of a major consequences of Bishop and Varmus's work was the demonstration that proto-oncogene activation plays a central role in human cancer, not just in virus-induced animal cancers. In the years following the discovery of \emph{c-src}, researchers identified activated oncogenes in a wide variety of human cancers, including \emph{RAS} mutations in pancreatic and colorectal cancers, \emph{MYC} amplification in breast and lung cancers, \emph{BCR-ABL} fusion in chronic myeloid leukemia, and \emph{HER2} amplification in breast cancer. The identification of these oncogenes has provided the basis for the development of targeted cancer therapies---drugs that specifically inhibit the activity of the oncogenic protein---that have transformed the treatment of many cancers.

\section{Impact on Modern Medicine}

The discovery that oncogenes arise from proto-oncogenes has had a transformative impact on modern medicine, particularly in the field of oncology. The identification of specific oncogenes in different types of cancer has enabled the development of targeted cancer therapies---drugs that selectively inhibit the activity of the oncogenic protein---that are more effective and less toxic than traditional chemotherapy.

The most striking example of this is the development of imatinib (Gleevec) for the treatment of chronic myeloid leukemia (CML). CML is caused by the chromosomal translocation t(9;22), which produces the BCR-ABL fusion protein---a constitutively active tyrosine kinase that drives uncontrolled cell proliferation. Imatinib, a small-molecule inhibitor of BCR-ABL, was developed based on the understanding of the molecular basis of CML that originated with Bishop and Varmus's work. Imatinib has transformed the treatment of CML from a uniformly fatal disease to a chronic condition with a near-normal life expectancy, and it is widely regarded as a notable success stories in the history of targeted cancer therapy.

The identification of HER2 amplification in breast cancer has led to the development of trastuzumab (Herceptin), a monoclonal antibody that targets the HER2 receptor and has significantly improved survival for patients with HER2-positive breast cancer. The identification of BRAF mutations in melanoma has led to the development of vemurafenib and dabrafenib, inhibitors of the mutant BRAF kinase that have dramatically improved survival for patients with metastatic melanoma. The identification of ALK rearrangements in non-small cell lung cancer has led to the development of crizotinib and other ALK inhibitors. These targeted therapies represent the direct clinical translation of Bishop and Varmus's basic research.

The concept that cancer is fundamentally a genetic disease---caused by mutations in specific genes that regulate cell growth and division---has also led to the development of cancer genomics and precision oncology. Large-scale projects such as The Cancer Genome Atlas (TCGA) and the International Cancer Genome Consortium (ICGC) have systematically cataloged the genetic mutations in thousands of cancers, and this information is now being used to guide the selection of targeted therapies for individual patients. Precision oncology---the matching of patients with the most appropriate targeted therapy based on the molecular profile of their tumor---is a direct legacy of Bishop and Varmus's discovery.

Beyond targeted therapy, the identification of proto-oncogenes and their mechanisms of activation has influenced the development of cancer diagnostics, prognostics, and prevention strategies. Molecular tests for specific oncogene mutations (such as \emph{RAS} mutations, \emph{BRAF} mutations, and \emph{HER2} amplification) are now routine in clinical oncology, and they are used to predict which patients will respond to specific therapies. The understanding of how environmental factors (such as tobacco smoke, UV radiation, and certain chemicals) activate proto-oncogenes has also informed cancer prevention strategies.

In the broader context of biology, the discovery that oncogenes arise from proto-oncogenes has provided insights into the fundamental mechanisms of cell growth and proliferation. Proto-oncogenes encode the normal signaling molecules that control cell growth, and their study has illuminated the complex networks of growth factor signaling, signal transduction, and gene expression that regulate the cell cycle. This knowledge has implications far beyond cancer, extending to developmental biology, regenerative medicine, and the understanding of growth disorders.

\section{Legacy}

J. Michael Bishop continued his work at the University of California, San Francisco, where he served as chancellor from 1998 to 2009. He has been a vocal advocate for basic research and for the translation of scientific discoveries into clinical practice. Bishop has written extensively about the process of scientific discovery and the challenges facing the scientific enterprise, and he remains an influential figure in American science.

Harold Varmus continued his work at UCSF and later became the director of the National Institutes of Health (NIH) from 1993 to 1999. During his tenure at the NIH, Varmus championed the Human Genome Project, the establishment of publicly accessible databases of scientific research, and the globalization of biomedical research. After leaving the NIH, Varmus became the director of the Memorial Sloan Kettering Cancer Center, where he continued to make contributions to cancer research and policy. He died in 2024.

The work of Bishop and Varmus demonstrates the power of basic scientific research to transform medicine. By pursuing a fundamental question about the origin of viral oncogenes, they discovered a principle that has revolutionized our understanding of cancer and has led to the development of life-saving targeted therapies. Their discovery that oncogenes arise from proto-oncogenes is one of the major findings in the history of cancer research, and it continues to inspire new approaches to the prevention, diagnosis, and treatment of cancer.


\section{Modern Developments}

Bishop and Varmus' proto-oncogene work has led to modern precision oncology. Understanding that cancer arises from mutated normal genes has led to targeted therapies (imatinib for BCR-ABL, vemurafenib for BRAF). Liquid biopsies detect circulating tumor DNA to monitor treatment response and detect resistance. Understanding of oncogene addiction has led to combination therapies that prevent resistance. CAR-T cell therapy targets cancer-specific antigens identified through oncogene research.


\section{Bibliography}

\begin{thebibliography}{9}
\bibitem{nobel-1989}
Nobel Prize Outreach, ``The Nobel Prize in Physiology or Medicine 1989,''\emph{NobelPrize.org}.\\
\url{https://www.nobelprize.org/prizes/medicine/1989/summary/}

\bibitem{lecture-1989}
J. Michael Bishop, ``Nobel Lecture,''\emph{NobelPrize.org}. Winner-authored primary source; downloadable from the official Nobel Prize archive.\\
\url{https://www.nobelprize.org/prizes/medicine/1989/bishop/lecture/}
\end{thebibliography}
